Mục tiêu tổng quát của đồ án là xây dựng và đánh giá CausClass, một khung phân tích video lớp học thống nhất từ nhận diện hành vi thị giác đến khám phá liên kết dự báo theo thời gian và sinh báo cáo quan sát dựa trên bằng chứng. Hệ thống được thiết kế để hỗ trợ người dùng trích xuất bằng chứng quan sát từ video, tóm tắt hành vi theo thời gian, phát hiện các liên kết dự báo giữa hành vi và trình bày kết quả theo cách có thể kiểm tra lại. Đầu ra của hệ thống là công cụ hỗ trợ phân tích, không thay thế đánh giá chuyên môn của giáo viên hoặc nhà nghiên cứu giáo dục.

Định hướng giải pháp thứ nhất nhằm xử lý hạn chế của nhận diện hành vi trong lớp học đông người. Đồ án sử dụng YOLO26n làm detector nền và bổ sung khối Multi-Head Self-Attention để tạo biến thể YOLO26n-MHSA. Mục tiêu của thay đổi này là kiểm tra liệu ngữ cảnh không gian xa giữa đầu, tay, bàn học và các vật thể nhỏ có giúp phân biệt những hành vi dễ nhầm lẫn như đọc, viết, cúi đầu hoặc sử dụng điện thoại hay không.

Định hướng giải pháp thứ hai nhằm chuyển đầu ra thị giác rời rạc thành dữ liệu thời gian có thể phân tích. Các phát hiện theo khung hình được liên kết bằng ByteTrack để tạo sự kiện hành vi có thời điểm bắt đầu--kết thúc, sau đó được gom vào các time bin cố định. Nhờ vậy, video không chỉ được biểu diễn bằng danh sách hộp phát hiện, mà còn bằng chuỗi thời gian đa biến, trong đó mỗi giá trị có thể truy ngược về các sự kiện và khung hình tương ứng.

Định hướng giải pháp thứ ba nhằm khám phá cấu trúc dự báo giữa các hành vi nhưng vẫn kiểm soát độ phức tạp của đồ thị. AERCA được dùng để tạo tín hiệu cấu trúc ban đầu và kiểm chứng các đồ thị ứng viên. LLM chỉ đóng vai trò đề xuất chỉnh sửa cấu trúc dựa trên tri thức miền và phản hồi từ các lần thử trước; mỗi đề xuất phải được đánh giá lại bằng masked AERCA và tiêu chí kết hợp lỗi dự báo với phạt số cạnh. Cách thiết kế này tận dụng tri thức ngôn ngữ nhưng không để LLM quyết định đồ thị cuối cùng.

Định hướng giải pháp thứ tư nhằm tạo báo cáo quan sát dễ đọc nhưng không vượt quá bằng chứng. Báo cáo phải mô tả dữ kiện phát hiện và thống kê thời gian trước, sau đó mới diễn giải các cạnh đồ thị như liên kết dự báo cần kiểm tra. Hệ thống không suy luận ý định học sinh, không đưa ra kết luận kỷ luật và không khẳng định chất lượng dạy học nếu không có bằng chứng ngoài video.

Phạm vi đánh giá của đồ án gồm ba phần. Phần thứ nhất đánh giá bộ phát hiện hành vi trên dữ liệu lớp học bằng các chỉ số như Precision, Recall và mAP. Phần thứ hai đánh giá tìm kiếm đồ thị trên 100 bộ kiểm thử thời gian tổng hợp có đồ thị chuẩn, nhằm đo Precision, Recall và F1 của các chiến lược tìm kiếm. Phần thứ ba minh họa toàn bộ quy trình trên video lớp học thật để kiểm tra khả năng truy vết từ video đến chuỗi thời gian, đồ thị và báo cáo. Các kết quả được diễn giải theo hướng bằng chứng quan sát và liên kết dự báo, không phải kết luận nhân quả chắc chắn.